\section*{Chapter \ref{ch:dv:multi-asset-options} --- Multi-Asset Options}

\begin{solution}{exo:dv:multi-asset-options:1}
$\sigma_B^2=0.25\times0.04+0.25\times0.09+2\times0.25\times0.5\times0.2\times0.3=0.0475$: 21.8\%, below the average volatility of
25\%.
\end{solution}

\begin{solution}{exo:dv:multi-asset-options:2}
A best-of call pays on the largest performance, which is far above the average when the assets move independently and close to
it when they move together: correlation shrinks its upside. A basket call pays on the average, whose volatility rises with
correlation: less diversification, more option value.
\end{solution}

\begin{solution}{exo:dv:multi-asset-options:3}
$\sum w_i^2\sigma_i^2=0.25(0.0625+0.1225)=0.04625$; $(\sum w_i\sigma_i)^2=0.09$, so the cross term is $0.04375$;
$\bar\rho=(0.0576-0.04625)/0.04375=0.26$.
\end{solution}

\begin{solution}{exo:dv:multi-asset-options:4}
Let $X$ be domestic per foreign. The domestic value of the foreign asset, $SX$, is a domestic traded asset, so under the domestic
measure it drifts at $r_d-q$. The exchange rate drifts at $r_d-r_f$. By Itô, the drift of $S=(SX)/X$ is
$(r_d-q)-(r_d-r_f)+\sigma_X^2-\operatorname{cov}(\ln SX,\ln X)/dt$, and
$\operatorname{cov}(\ln SX,\ln X)/dt=\rho\sigma_S\sigma_X+\sigma_X^2$. The drift is $r_f-q-\rho\sigma_S\sigma_X$, so the quanto
forward is the foreign forward times $e^{-\rho\sigma_S\sigma_XT}$.
\end{solution}

\begin{solution}{exo:dv:multi-asset-options:5}
The note embeds a worst-of put sold by the client (the capital-at-risk leg) and a worst-of digital coupon bought by the client;
the dealer is long the put and short the digital, and both lose value for the dealer as correlation rises. Index options, or
index variance against single-stock variance (a reverse dispersion), carry the opposite correlation exposure and offset part of
it.
\end{solution}

\begin{solution}{exo:dv:multi-asset-options:6}
About $+28$: the hedge amounts cancel the first-order exposure to each member's variance, and what is left is a second-order
term, small against the 270\,600 that a fifteen-point correlation gap produces.
\end{solution}

\begin{solution}{exo:dv:multi-asset-options:7}
0.736, between the constant-correlation prices at 0.6 (0.707) and 0.8 (0.748), and well above the 0.673 at the index's
at-the-money \omterm{def:dv:multi-asset-options:implied}{implied correlation} of about 0.4: under \omterm{def:dv:multi-asset-options:local}{local correlation} the shares move together in exactly the falling
scenarios that decide the digital.
\end{solution}

\begin{solution}{exo:dv:multi-asset-options:8}
It is a premium, not a carry: the book is short correlation and short index volatility, and it loses heavily in a sell-off,
when correlation rises towards one and index volatility spikes. The published evidence also says that realistic trading
frictions make the premium hard to capture. The average gap is compensation for that risk.
\end{solution}

\begin{solution}{pb:dv:multi-asset-options:1}
\textbf{1.} $\sum w_i^2\sigma_i^2=0.00430$ (43.0 variance points); $C=0.0798$ (798).
\textbf{2.} 6.6\% at zero correlation; 29.0\% at one, the weighted average volatility.
\textbf{3.} Diversification: with correlations below one, members' moves partly cancel.
\textbf{4.} 21.9\%, 19.2\% and 17.0\%.
\textbf{5.} 19.3\% at both: a flat index smile.
\textbf{6.} 0.549, 0.408 and 0.308.
\textbf{7.} 0.559, from the index's strip volatility of 22.1\%.
\textbf{8.} Index options below the money price the members falling together, which is a higher correlation.
\textbf{9.} $\rho(x)=0.384-3.62x+16.3x^2$, clipped to $[0,1]$.
\textbf{10.} 0.91 at $x=-0.1$ and 0.19 at $x=+0.1$.
\textbf{11.} $100\,000/(2\times22.12)=2\,261$ per variance point.
\textbf{12.} $N\bigl(w_i^2(1-\bar\rho)+\bar\rho w_i\sum_jw_j\sigma_j/\sigma_i\bigr)$: the index variance's sensitivity to each member's
variance at the \omterm{def:dv:multi-asset-options:implied}{implied correlation}; they cancel the first-order exposure to each member's volatility.
\textbf{13.} $2\,261\times0.15\times798=270\,600$.
\textbf{14.} 271\,000, 34\,000, from 215\,000 to 326\,000.
\textbf{15.} The realised average correlation, 0.409, minus 0.559: $-0.15$ per unit notional to the long side, $+0.15$ to the
short.
\textbf{16.} When correlation realises above implied, typically in a sell-off; the loss is $N(\hat\rho-\bar\rho)C$ plus the
residual from volatility moves, and correlation can go to one: about $2\,261\times0.441\times798\approx800\,000$ from correlation
alone if members realise their implied volatilities.
\textbf{17.} Implied has averaged above realised: 39.5\% against 32.5\% for the S\&P~500, 46.0\% against 35.5\% for the Dow Jones 30,
in Driessen, Maenhout and Vilkov's sample.
\textbf{18.} Investors pay to hedge correlation spikes (they are long index puts), and the trade that sells it needs many option
positions and bears crash risk; frictions limit arbitrage.
\textbf{19.} \omterm{def:dv:multi-asset-options:implied}{Implied correlation} 0.549 at 90, 0.408 at 100 and 0.308 at 110 (0.559 in variance terms); the \omterm{def:dv:multi-asset-options:dispersion}{dispersion trade} makes
270\,600 on a \omterm{def:dv:variance-swaps-and-volatility-derivatives:veganot}{vega notional} of 100\,000 when realised correlation is fifteen points below implied (simulated 271\,000 $\pm$ 34\,000).
\textbf{20.} The joint behaviour of its members, above all how their correlation changes when the market falls.
\end{solution}

\begin{solution}{iq:dv:multi-asset-options:1}
A basket call's value rises with correlation (less diversification, more volatility of the average); a worst-of put and a best-of
call fall with it (the extremes of the performances move closer to the average as correlation rises).
\iqlookfor{the three signs and the mechanism.}
\end{solution}

\begin{solution}{iq:dv:multi-asset-options:2}
The single correlation that reconciles the index's implied variance with its members': $(\sigma_I^2-\sum w_i^2\sigma_i^2)/(\text{cross
term})$. Index skews are steeper than single-stock skews, so the index's volatility at low strikes implies a higher correlation: the
market prices stocks falling together.
\iqlookfor{the formula and the \omterm{def:dv:multi-asset-options:skew}{correlation skew}.}
\end{solution}

\begin{solution}{iq:dv:multi-asset-options:3}
Short index variance (or straddles), long members' variance in amounts that cancel each member's volatility exposure at the \omterm{def:dv:multi-asset-options:implied}{implied
correlation}. The residual is short correlation: P\&L $\approx N(\bar\rho-\hat\rho)C$. Exposed to correlation spikes in sell-offs,
to the volatility of the weights and hedge ratios, and to transaction costs across many legs.
\iqlookfor{the legs, the weights and the residual exposure.}
\end{solution}

\begin{solution}{iq:dv:multi-asset-options:4}
Under the domestic measure the foreign asset's drift is $r_f-q-\rho\sigma_S\sigma_X$, with $X$ in domestic per foreign units (Girsanov
on the change of numeraire). For a yen index paid in dollars, $X$ is dollars per yen; if the index falls when the yen strengthens,
$\rho<0$, and the quanto forward is above the yen forward.
\iqlookfor{the drift and the sign convention.}
\end{solution}

\begin{solution}{iq:dv:multi-asset-options:5}
The matrix is not positive definite: pairwise estimates from different windows or data, or hand-set entries, are inconsistent. Project
it onto the nearest correlation matrix (clip negative eigenvalues and rescale the diagonal, or an alternating-projections method), or
build it from a factor model, and check the change in prices.
\iqlookfor{positive definiteness and a principled repair.}
\end{solution}

\begin{solution}{iq:dv:multi-asset-options:6}
Measure: the book's value change for a shift of all pairwise correlations (and for a \omterm{def:dv:multi-asset-options:skew}{correlation skew} shift), by bucket of maturity
and moneyness of the worst performer. Hedge: with index options or index variance against single-stock variance, which carry
correlation exposure of the opposite sign, knowing the hedge is imperfect (index correlation is not the basket's). Reserve: for the
unhedged correlation at a stressed level and for the model (local against constant correlation).
\iqlookfor{a correlation sensitivity, a proxy hedge and a reserve.}
\end{solution}
